% =====================================================================
%  Chapter 7 — Evaluation: Transformation Experiments
%  Drop-in for the existing thesis. Sections 7.2–7.6.
%  Uses booktabs and siunitx-style numbers; falls back gracefully if
%  those packages are unavailable. No external graphics.
% =====================================================================

\section{Experimental Setup}\label{sec:transformation-setup}

The transformation experiments operationalise the second evaluation
dimension introduced in Section~3.4: they assess to what extent the
combination of \textsc{AgentOscin} and the OSCIN-based pipeline can
\emph{round-trip} an agentic AI system through a framework-agnostic
semantic representation. Three complementary pipelines are exercised
over an identical example corpus:

\begin{description}
  \item[Extraction] (\textit{source}~$\rightarrow$~\textsc{ttl}): the
    parser--populator chain ingests a framework-specific code base and
    materialises an \textsc{AgentOscin} A-Box. Used to quantify how
    much of an existing system survives the lift into the ontology.
  \item[Generation] (\textsc{ttl}~$\rightarrow$~\textit{source}): a
    target-framework code generator emits source code from a populated
    A-Box. Used to verify that the ontology contains enough information
    to reconstitute a runnable artefact.
  \item[Round-trip] (\textit{source}~$\rightarrow$~\textsc{ttl}~$\rightarrow$~\textit{source$'$}~$\rightarrow$~\textsc{ttl$'$}):
    the two halves are composed. In the
    \emph{same-framework} mode the two ends share a language and we can
    additionally compare source trees and execution traces; in the
    \emph{cross-framework} mode the source language differs, so the
    comparison is restricted to the two A-Boxes and (where applicable)
    canonical mapping rules.
\end{description}

\paragraph{Corpus.} The corpus consists of 21 publicly available
example systems collected from the official documentation, tutorials
and reference repositories of the three target frameworks
(Table~\ref{tab:eval-corpus}). The selection deliberately spans
distinct architectural patterns — sequential crews, event-driven flows,
shared-state graphs, group chats — so that the metrics exercise
different parts of the ontology rather than only the agent--task--tool
core. Three of the AutoGen examples (\texttt{company-research},
\texttt{literature-review}, \texttt{travel-planning}) are distributed
as Jupyter notebooks; the AutoGen parser concatenates code cells via
\texttt{nbformat} before feeding them to the AST extractor, so they
participate in the extraction benchmark on equal footing. They are
excluded from the round-trip and generation tables because those
pipelines assume a \texttt{.py} entry point and the generators only
emit Python modules, not notebooks; widening the round-trip pipeline
to round-trip notebooks would require a notebook emitter and is left
for future work.

\begin{table}[ht]
  \centering
  \small
  \caption{Example corpus used in the transformation experiments.
           The three notebook-only AutoGen examples participate in
           extraction but not in round-trip / generation.}
  \label{tab:eval-corpus}
  \resizebox{\textwidth}{!}{%
  \begin{tabular}{l c c l}
    \toprule
    Framework & \# extraction & \# round-trip & Architectural patterns covered \\
    \midrule
    CrewAI    & 7 & 7 & sequential crew, hierarchical crew, \texttt{@Flow} state machine \\
    LangGraph & 6 & 6 & ReAct, plan-and-execute, RAG, reflexion, memory-augmented \\
    AutoGen   & 8 & 5 & \texttt{RoundRobinGroupChat}, \texttt{SelectorGroupChat}, code review, data analysis, +3 notebooks \\
    \midrule
    Total     & 21 & 18 & \\
    \bottomrule
  \end{tabular}
  }
\end{table}

\paragraph{Metric catalogue.}
The benchmark harness in \texttt{evaluation/} composes seven metric
modules, each implementing a uniform
\textsc{compute}/\textsc{row}/\textsc{summarize} contract.
Table~\ref{tab:metric-catalogue} summarises which metric is reported by
which pipeline.

\begin{table}[ht]
  \centering
  \small
  \caption{Metric catalogue and their use across pipelines.
           \textsc{ttl} denotes the populated A-Box;
           same-fw and cross-fw refer to the two round-trip modes.}
  \label{tab:metric-catalogue}
  \resizebox{\textwidth}{!}{%
  \begin{tabular}{l p{6.4cm} c c c c}
    \toprule
                                & & \multicolumn{4}{c}{Pipeline}\\
    \cmidrule(lr){3-6}
    Metric                      & What is compared
                                & Extr. & Gen. & RT same & RT cross\\
    \midrule
    \texttt{ttl\_intrinsic}     & counts/density of one A-Box
                                & \checkmark & \checkmark & & \\
    \texttt{ttl\_pairwise}      & two A-Boxes: individual / property /
                                  triple precision-recall-F$_1$,
                                  literal overlap
                                & (\checkmark) & (\checkmark) &
                                  \checkmark & \checkmark\\
    \texttt{ttl\_fuzzy\_match}  & two A-Boxes with name-tolerant
                                  entity alignment
                                & (\checkmark) & & \checkmark &
                                  \checkmark\\
    \texttt{ast\_diff}          & two source trees in the same
                                  language: imports, classes,
                                  decorators, state fields, graph
                                  calls
                                & & (\checkmark) & \checkmark & \\
    \texttt{syntax\_validity}   & every generated \texttt{.py}:
                                  parses with \texttt{ast.parse} and
                                  imports resolve
                                & & \checkmark & & \\
    \texttt{execution\_trace}   & subprocess run of the original and
                                  the generated entry point
                                & & (\checkmark) & \checkmark & \\
    \texttt{mapping\_conformance}& canonical CrewAI~Flow~$\leftrightarrow$~LangGraph
                                  rule hit-rate
                                & & & & (\checkmark)\\
    \bottomrule
    \multicolumn{6}{l}{\footnotesize\checkmark{} = always reported;
      (\checkmark) = reported when the optional input
      (ground-truth, reference source, re-extraction) is available.}
  \end{tabular}
  }
\end{table}

For pairwise comparisons of two graphs $G_{\text{ref}}$ and
$G_{\text{cand}}$ the harness reports precision, recall and $F_1$ at
three granularities, computed with class-aware local-name normalisation
so that differing namespace prefixes do not penalise the score:
\begin{align*}
  P_{\textsc{lvl}} &= \frac{|R_{\textsc{lvl}}\cap C_{\textsc{lvl}}|}{|C_{\textsc{lvl}}|}, \quad
  R_{\textsc{lvl}} = \frac{|R_{\textsc{lvl}}\cap C_{\textsc{lvl}}|}{|R_{\textsc{lvl}}|}, \quad
  F_{1,\textsc{lvl}} = \frac{2P_{\textsc{lvl}}R_{\textsc{lvl}}}{P_{\textsc{lvl}}+R_{\textsc{lvl}}},\\
  \text{level} &\in \{\text{individual}, \text{property}, \text{triple}\}.
\end{align*}
The literal-overlap signal is a Jaccard coefficient over typed literal
strings. The fuzzy-match metric relaxes individual-level identity by
greedily aligning entities across the two graphs using a similarity
score that combines \texttt{difflib.SequenceMatcher} ratio with a
token-set Jaccard bonus and accepts any pair scoring above
$\tau{=}0{.}55$; we report the number of aligned pairs and the average
score. The structural \texttt{ast\_diff} metric extracts nine feature
sets per source tree (imports, classes, class bases, function names,
decorators, decorator-argument literals, \texttt{TypedDict}/Pydantic
state fields, state annotations, LangGraph \texttt{StateGraph} method
calls), computes per-feature $F_1$, and reports the macro mean as
\textsc{ast\_overall\_f1}. The \texttt{execution\_trace} metric runs
both entry points in isolated subprocesses with
\texttt{dotenv.load\_dotenv()} and compares exit status, runtime ratio
and a token-level Jaccard over standard-output streams.

\paragraph{Reproducibility.}
All numbers reported in the following sections were produced by the
checked-in harness and can be regenerated with three deterministic
commands:
\begin{verbatim}
python -m evaluation.benchmark --pipeline extraction
python -m evaluation.benchmark --pipeline generation \
       --target-framework <fw> --from-extraction
python -m evaluation.benchmark --pipeline roundtrip \
       --skip-execution [--target-framework <fw>]
\end{verbatim}
Per-example reports are written under \texttt{output/<pipeline>/} as
\texttt{report.md}, \texttt{report.json} and an aggregated
\texttt{benchmark.csv}. The CSVs underlying every table below are
included in the supplementary material.

% =====================================================================
\section{Extraction Evaluation}\label{sec:extraction-eval}
% =====================================================================

The extraction pipeline lifts each of the 21 example systems
(including the three Jupyter-notebook AutoGen examples; see
Table~\ref{tab:eval-corpus}) into an \textsc{AgentOscin} A-Box. Two
complementary measurements are reported. First, an \emph{intrinsic}
audit (Table~\ref{tab:extraction-intrinsic}) on the full 21-system
corpus, recording whether the produced graphs are non-trivial in size
and use the ontology breadth. Second, a \emph{gold-anchored}
comparison (Section~\ref{sec:llm-baseline}) on a 6-system subset for
which a reference A-Box was hand-authored from source according to
the annotation guideline in
\texttt{docs/annotation\_guideline.md}; against that reference both
the deterministic pipeline and an LLM extractor that uses the same
ontology and prompt are scored. All 21 systems extracted cleanly (no
parser crash, no populator validation error).

\begin{table}[ht]
  \centering
  \small
  \caption{Intrinsic A-Box statistics per extraction. Triples is
    the total number of triples written, Indiv. counts named
    individuals, Prop. counts distinct ontology properties used and
    Density is the number of A-Box triples per individual.
    Rows marked $\dagger$ are extracted from Jupyter notebooks
    via \texttt{nbformat}.}
  \label{tab:extraction-intrinsic}
  \resizebox{\textwidth}{!}{%
  \begin{tabular}{l l r r r r}
    \toprule
    Framework & Example & Triples & Indiv. & Prop. & Density\\
    \midrule
    CrewAI    & academic-research-flow & 671 & 45 & 54 & 4.18\\
    CrewAI    & code-review            & 678 & 43 & 44 & 4.54\\
    CrewAI    & comprehensive          & 657 & 42 & 56 & 4.14\\
    CrewAI    & content-pipeline       & 658 & 36 & 42 & 4.86\\
    CrewAI    & email-flow             & 670 & 39 & 44 & 4.80\\
    CrewAI    & self-eval-loop-flow    & 634 & 37 & 40 & 4.08\\
    CrewAI    & tech-blog              & 623 & 27 & 35 & 5.19\\
    \addlinespace
    LangGraph & ReAct                  & 571 & 20 & 34 & 4.40\\
    LangGraph & drafter                & 563 & 18 & 33 & 4.44\\
    LangGraph & memoryagent            & 553 & 18 & 28 & 3.89\\
    LangGraph & ragagent               & 602 & 27 & 38 & 4.41\\
    LangGraph & research-assistant     & 610 & 27 & 39 & 4.70\\
    LangGraph & tech-blog              & 621 & 29 & 31 & 4.76\\
    \addlinespace
    AutoGen   & code-review            & 593 & 26 & 29 & 4.23\\
    AutoGen   & company-research$^\dagger$ & 610 & 28 & 29 & 4.54\\
    AutoGen   & comprehensive          & 586 & 23 & 29 & 4.48\\
    AutoGen   & content-pipeline       & 599 & 27 & 29 & 4.30\\
    AutoGen   & data-analysis-0.4      & 595 & 27 & 28 & 4.15\\
    AutoGen   & literature-review$^\dagger$ & 609 & 27 & 29 & 4.67\\
    AutoGen   & tech-blog              & 584 & 24 & 26 & 4.21\\
    AutoGen   & travel-planning$^\dagger$ & 613 & 28 & 26 & 4.64\\
    \midrule
    \multicolumn{2}{l}{Mean (CrewAI, $n{=}7$)}    & 655.9 & 38.4 & 45.0 & 4.54\\
    \multicolumn{2}{l}{Mean (LangGraph, $n{=}6$)} & 586.7 & 23.2 & 33.8 & 4.43\\
    \multicolumn{2}{l}{Mean (AutoGen, $n{=}8$)}   & 598.6 & 26.3 & 28.1 & 4.40\\
    \bottomrule
  \end{tabular}
  }
\end{table}

Three observations follow directly from
Table~\ref{tab:extraction-intrinsic}. First, the per-system A-Box size
varies by less than $\pm$10\,\% within a framework — including the
three notebook-sourced AutoGen examples, which fall in the same
range as the \texttt{.py} ones — indicating that extraction quality
is dominated by the static fixture written into the ontology by the
populator (system metadata, namespaces, \texttt{rdfs:label}s) rather
than by example-specific noise or by the input format. Second, CrewAI examples consume
the broadest property vocabulary (mean 45 properties) because the
\texttt{@CrewBase}/\texttt{@Flow} stack lets the parser observe agent
goals, backstories, structured task outputs and explicit step
sequencing in the same source tree; LangGraph examples instead rely on
free-form node functions and an \texttt{Annotated[\dots]} state
schema, which the parser models with fewer properties (mean 33.8) but
still produces a comparable number of total triples through the
fine-grained \texttt{WorkflowStep} graph. Third, the information
density (4.0–5.2 triples per individual across all frameworks)
indicates that the populator emits a stable amount of metadata per
extracted entity rather than degenerating into a sparse type-only
graph; this is the precondition for the pairwise round-trip metrics
to be meaningful in the next section.

% =====================================================================
\section{Generation Evaluation}\label{sec:generation-eval}
% =====================================================================

The generation pipeline is the inverse of extraction: the populated
A-Box is read by \texttt{oscin.reader} into the intermediate
representation and emitted as target-framework source by the
generators in \texttt{oscin/generators/}. To stress the generators
with realistic, structurally diverse inputs we feed the 18 A-Boxes
produced in the extraction step into each of the three generators in
turn ($n{=}54$ generation runs). The two static metrics reported are
\emph{syntax rate} (the fraction of generated \texttt{.py} files that
parse with \texttt{ast.parse}) and \emph{import rate} (the fraction of
generated files for which every top-level import resolves against the
current Python environment).

\begin{table}[ht]
  \centering
  \small
  \caption{Generation pipeline syntactic validity, aggregated over
           the 18-system corpus. Range is per-system min/max of the
           import rate.}
  \label{tab:generation-validity}
  \resizebox{\textwidth}{!}{%
  \begin{tabular}{l c c c c}
    \toprule
    Target framework & Runs & Mean syntax rate & Mean import rate & Range\\
    \midrule
    CrewAI     & 18 & 1.000 & 0.721 & [0.333, 1.000]\\
    LangGraph  & 18 & 1.000 & 1.000 & [1.000, 1.000]\\
    AutoGen    & 18 & 1.000 & 1.000 & [1.000, 1.000]\\
    \bottomrule
  \end{tabular}
  }
\end{table}

All 54 generation runs produce code that parses without
\texttt{SyntaxError}, which establishes a structural-correctness floor:
the templated emitters in \texttt{oscin/generators/} are not fragile
under unfamiliar A-Boxes. The LangGraph and AutoGen targets reach a
perfect import resolution rate because their generators consolidate the
generated module into a single \texttt{main.py} that depends only on
the framework runtime and \texttt{python-dotenv}, both of which are
installed in the evaluation environment. The CrewAI generator emits a
multi-file project — \texttt{config/agents.yaml},
\texttt{config/tasks.yaml}, \texttt{crew.py}, \texttt{main.py} — that
in some cases imports tool stubs or local modules whose names are
borrowed from the source A-Box rather than synthesised; these stubs
are flagged as \emph{external} in the populator and intentionally
preserved as references. The 0.721 mean import rate (range
0.333–1.000) reflects this: in 5 of 18 cases at least one such
external reference cannot be resolved against PyPI, which is a
modelling choice rather than a generation failure (Section
\ref{sec:transformation-discussion}).

% =====================================================================
\section{Round-Trip and Cross-Framework Evaluation}\label{sec:roundtrip-eval}
% =====================================================================

The round-trip pipeline composes extraction and generation:
\[
  S \xrightarrow{\text{extract}_{\,f}} G_1
    \xrightarrow{\text{generate}_{\,t}} S'
    \xrightarrow{\text{extract}_{\,t}} G_2,
\]
where the source framework $f$ and target framework $t$ either
coincide (\emph{same-framework} mode) or differ
(\emph{cross-framework} mode). The two A-Boxes $G_1, G_2$ are then
compared with \texttt{ttl\_pairwise} and the name-tolerant
\texttt{ttl\_fuzzy\_match}; in same-framework mode the source trees
$S, S'$ are additionally compared with \texttt{ast\_diff}.

\subsection{Same-framework round-trip}

The same-framework round-trip is the strictest test of the
methodology: a system passes through both directions of the pipeline
in the language it was authored in. Drift in either direction surfaces
here.

\begin{table}[ht]
  \centering
  \small
  \caption{Same-framework round-trip on the 18-system corpus
           ($G_1$ vs.\ $G_2$ for the A-Box columns; original source
           vs.\ regenerated source for \textsc{ast}~F$_1$).}
  \label{tab:roundtrip-same}
  \resizebox{\textwidth}{!}{%
  \begin{tabular}{l l c c c c c}
    \toprule
    Framework & Example
              & Triple F$_1$ & Property F$_1$ & Individual F$_1$
              & Fuzzy avg. & \textsc{ast}~F$_1$\\
    \midrule
    CrewAI    & academic-research-flow & 0.877 & 0.971 & 0.966 & 1.000 & 0.742\\
    CrewAI    & code-review            & 0.910 & 1.000 & 1.000 & 1.000 & 0.873\\
    CrewAI    & comprehensive          & 0.880 & 0.936 & 0.976 & 1.000 & 0.794\\
    CrewAI    & content-pipeline       & 0.909 & 1.000 & 1.000 & 1.000 & 0.886\\
    CrewAI    & email-flow             & 0.925 & 1.000 & 1.000 & 1.000 & 0.789\\
    CrewAI    & self-eval-loop-flow    & 0.910 & 1.000 & 1.000 & 1.000 & 0.924\\
    CrewAI    & tech-blog              & 0.914 & 1.000 & 1.000 & 1.000 & 0.972\\
    \addlinespace
    LangGraph & ReAct                  & 0.758 & 0.971 & 0.952 & 1.000 & 0.696\\
    LangGraph & drafter                & 0.739 & 0.955 & 0.947 & 1.000 & 0.652\\
    LangGraph & memoryagent            & 0.817 & 0.966 & 1.000 & 1.000 & 0.739\\
    LangGraph & ragagent               & 0.861 & 0.974 & 0.981 & 1.000 & 0.591\\
    LangGraph & research-assistant     & 0.874 & 1.000 & 1.000 & 1.000 & 0.581\\
    LangGraph & tech-blog              & 0.863 & 0.969 & 0.951 & 1.000 & 0.692\\
    \addlinespace
    AutoGen   & code-review            & 0.862 & 1.000 & 0.980 & 1.000 & 0.992\\
    AutoGen   & comprehensive          & 0.610 & 0.912 & 0.857 & 0.983 & 0.870\\
    AutoGen   & content-pipeline       & 0.861 & 1.000 & 0.981 & 1.000 & 0.986\\
    AutoGen   & data-analysis-0.4      & 0.816 & 0.982 & 0.943 & 1.000 & 0.970\\
    AutoGen   & tech-blog              & 0.871 & 1.000 & 1.000 & 1.000 & 0.989\\
    \midrule
    \multicolumn{2}{l}{Mean (CrewAI)}    & 0.904 & 0.987 & 0.992 & 1.000 & 0.854\\
    \multicolumn{2}{l}{Mean (LangGraph)} & 0.819 & 0.973 & 0.972 & 1.000 & 0.659\\
    \multicolumn{2}{l}{Mean (AutoGen)}   & 0.804 & 0.979 & 0.952 & 0.997 & 0.961\\
    \multicolumn{2}{l}{\textbf{Mean (overall)}} & \textbf{0.847} & \textbf{0.980} & \textbf{0.974} & \textbf{0.999} & \textbf{0.818}\\
    \bottomrule
  \end{tabular}
  }
\end{table}

Three regularities emerge from Table~\ref{tab:roundtrip-same}.

\textbf{Property and individual identity are essentially preserved.}
The mean property $F_1$ is 0.980 and the mean individual $F_1$ is
0.974; the fuzzy-alignment score is at ceiling (0.999). This means
that the \emph{set of ontology constructs in play} and the
\emph{set of named entities} survive the round-trip almost exactly,
i.e., the populator and the generators agree on the ontology
vocabulary they consume and produce.

\textbf{Triple-level $F_1$ is bounded below 1 by inherent
abstraction loss.} The mean triple $F_1$ of 0.847 is the headline
number for ontology idempotence. The 15-point gap to perfect
preservation is dominated by literal-valued properties whose source
form is paraphrased or canonicalised by the generator
(e.g., \texttt{hasTitle} regenerated as a slugified system name) and
by ordering predicates on \texttt{WorkflowStep} chains that the
re-extractor reconstructs in a different but semantically equivalent
order. The ontology idempotence is therefore approximate at the
triple level but \emph{structurally} faithful — a distinction that the
property/individual scores quantify.

\textbf{Source-level reconstruction quality follows the framework.}
The \textsc{ast}~F$_1$ score reaches 0.961 mean for AutoGen, where the
generator template tracks the source pattern almost line-for-line, but
drops to 0.659 for LangGraph because the generator emits a
canonicalised \texttt{StateGraph} skeleton with stable node names and
loses the ad-hoc helpers and prompt fragments that the original code
embedded. CrewAI sits in between (0.854) because its YAML-driven
configuration absorbs most of the structural variance into a stable
file layout. This is consistent with the role of the ontology as a
\emph{semantic} rather than \emph{syntactic} contract: it preserves
agents, tasks, tools, teams and workflow topology, but does not (and
should not) claim to preserve every implementation detail of the
original Python module.

\subsection{Competency-question survivability}\label{sec:cq-survivability}

The pairwise metrics in Table~\ref{tab:roundtrip-same} report on the
\emph{triple inventory} of $G_1$ and $G_2$ but cannot, on their own,
tell whether the questions the ontology was designed to answer
(Section~7.1, Table~7.2) are still answerable on the round-tripped
graph. To address this directly we encoded four representative
competency questions as SPARQL \texttt{SELECT} queries —
\textbf{CQ-1} (agents and their identifiers), \textbf{CQ-10} (team
coordination pattern), \textbf{CQ-14} (ordered workflow steps) and
\textbf{CQ-18} (tools and their input schemas) — chosen one from each
of the four broad concept clusters in the requirements catalogue
(Agent~Definition, Coordination, Perception/Workflow, Tools), and
evaluated each against $G_1$ and $G_2$ for every same-framework
round-trip system. A query result is recorded as \emph{preserved
exactly} iff the answer multiset on $G_2$ equals the answer multiset
on $G_1$ after local-name canonicalisation; we additionally report the
mean Jaccard coefficient over the union of the two answer sets to give
a graded view.

\begin{table}[ht]
  \centering
  \small
  \caption{Competency-question survivability across the 18-system
           same-framework round-trip. \emph{Exact} is the fraction
           of systems in which the SPARQL answer multiset on $G_2$
           equals the multiset on $G_1$ after local-name
           canonicalisation; \emph{Mean Jaccard} is the average
           overlap over the union of answer sets.}
  \label{tab:cq-survivability}
  \resizebox{\textwidth}{!}{%
  \begin{tabular}{l p{6.6cm} c c}
    \toprule
    CQ & Question & Exact & Mean Jaccard\\
    \midrule
    CQ-1  & Which agents are defined and how is each identified?       & 1.00 & 1.000\\
    CQ-10 & What coordination pattern does each team follow?           & 0.83 & 0.833\\
    CQ-14 & What is the step-by-step workflow structure (ordered)?     & 0.83 & 0.833\\
    CQ-18 & What tools are available and what input schema do they take? & 0.78 & 0.848\\
    \bottomrule
  \end{tabular}
  }
\end{table}

Three observations follow.
First, agent-identity questions (CQ-1) are preserved on \emph{every}
system; this corroborates the ceiling individual-$F_1$ score in
Table~\ref{tab:roundtrip-same} but adds the stronger guarantee that
the \texttt{?agent~ao:agentID~?id} pattern returns the same bindings
on $G_2$ as on $G_1$, not merely the same set of subjects.
Second, the coordination-pattern question (CQ-10) and the
workflow-structure question (CQ-14) are preserved on 15 of 18 systems;
the three losses come from systems where the generator emits a
canonical \texttt{Sequential} pattern in place of a more specific one
(\texttt{ReActLoop}, \texttt{ReflectionLoop}) when no reasoning-pattern
binding survived the A-Box. This is a recoverable gap — the
information is present in $G_1$ — and is the same root cause that
drives the LangGraph \texttt{ast\_diff} dip discussed above.
Third, the tool-schema question (CQ-18) is preserved exactly on
14 of 18 systems but with a higher mean Jaccard (0.848): tool
\emph{identities} are intact in every case but the
\texttt{ao:hasInputSchema} literal is paraphrased by the generator
in four systems, which is consistent with the literal-paraphrase
hypothesis in the previous subsection. We therefore conclude that the
round-trip preserves the answers the ontology was designed to support
\emph{exactly} for $\geq 0.78$ of systems on every CQ, and
\emph{semantically} (Jaccard $\geq 0.83$) on every system. The full
per-system breakdown is provided in
\texttt{output/cq\_survivability/summary.md}.

\subsection{End-to-end walked-through example}\label{sec:walkthrough}

To make the round-trip concrete, this subsection traces a single
system — the CrewAI \texttt{email-flow} example — through the three
stages of the pipeline. The example is representative rather than
best-case (Table~\ref{tab:roundtrip-same}: triple~$F_1$ 0.925,
property~$F_1$ 1.000, individual~$F_1$ 1.000,
\textsc{ast}~F$_1$ 0.789).

\paragraph{Stage 1 — original source.}
The system is authored as an \texttt{@CrewBase}-decorated crew with
three agents whose roles, goals and backstories are configured in a
YAML sidecar:

\begin{verbatim}
# crews/config/agents.yaml
email_filter_agent:
  role: >
    Senior Email Analyst
  goal: >
    Filter out non-essential emails like newsletters and
    promotional content, and identify important
    action-required emails that need immediate attention.
  backstory: >
    With years of experience in email management and
    communication triage, you are adept at quickly identifying
    important emails that require immediate action ...
\end{verbatim}

The orchestration around the crew is encoded as a CrewAI
\texttt{Flow} subclass with two \texttt{@start}/\texttt{@listen}
methods (excerpt):

\begin{verbatim}
class EmailAutoResponderFlow(Flow[AutoResponderState]):
    @start("wait_next_run")
    def fetch_new_emails(self):
        new_emails, updated = check_email(...)
        self.state.emails = new_emails

    @listen(fetch_new_emails)
    def generate_draft_responses(self):
        if len(self.state.emails) > 0:
            EmailFilterCrew().crew().kickoff(
                inputs={"emails": format_emails(self.state.emails)})
\end{verbatim}

\paragraph{Stage 2 — extracted A-Box ($G_1$).}
The CrewAI parser surfaces the YAML- and decorator-derived facts and
the populator turns them into 175 ABox triples (excerpt, namespaces
elided):

\begin{verbatim}
ex:email_flow a ao:AgenticSystem ;
    ao:containsAgent ex:Agent_email_filter_agent,
                     ex:Agent_email_action_agent,
                     ex:Agent_email_response_writer ;
    ao:containsOrchestration ex:Orchestration_EmailAutoResponderFlow ;
    ao:containsTeam ex:Team_EmailFilterCrew ;
    ao:hasSourceFramework "CrewAI" .

ex:Agent_email_filter_agent a ao:LLMAgent ;
    ao:agentID    "email_filter_agent" ;
    ao:agentRole  "Senior Email Analyst" ;
    ao:agentPrompt ex:AgentPrompt_email_filter_agent ;
    ao:hasGoal     ex:Goal_email_filter_agent .

ex:AgentPrompt_email_filter_agent a ao:Prompt ;
    ao:hasDirectiveFunction "DualDirective" ;
    ao:promptInstruction
        "Senior Email Analyst: Filter out non-essential emails ..." ;
    ao:promptContext
        "With years of experience in email management ..." .

ex:FlowStep_fetch_new_emails a ao:StartStep, ao:WorkflowStep ;
    ao:hasDecoratorArgument "wait_next_run" ;
    ao:nextStep ex:FlowStep_generate_draft_responses ;
    ao:stepOrder 1 .
\end{verbatim}

Three things to note. The agent's \texttt{role} string is preserved
verbatim as an \texttt{ao:agentRole} literal; the \texttt{goal} and
\texttt{backstory} are split across an \texttt{ao:Goal} individual and
the \texttt{promptInstruction}/\texttt{promptContext} of the prompt
individual respectively, materialising the
instruction--context--output decomposition that the ontology
prescribes. The \texttt{@start}/\texttt{@listen} pair becomes a
two-step \texttt{WorkflowPattern} with an explicit \texttt{stepOrder}
and a \texttt{nextStep} link, which is what allows CQ-14 above to
recover an ordered workflow purely from the A-Box.

\paragraph{Stage 3 — regenerated source.}
Feeding $G_1$ back through the CrewAI generator produces a five-file
project (\texttt{main.py}, \texttt{crews/email\_filter\_crew/email\_filter\_crew.py},
\texttt{config/agents.yaml}, \texttt{config/tasks.yaml},
\texttt{tools/create\_draft.py}). The regenerated YAML mirrors the
A-Box almost verbatim:

\begin{verbatim}
# generated/crews/email_filter_crew/config/agents.yaml
email_filter_agent:
  role: Senior Email Analyst
  goal: Filter out non-essential emails like newsletters
        and promotional content, and identify important
        action-required emails that need immediate attention.
  backstory: With years of experience in email management
        and communication triage, you are adept at quickly
        identifying important emails ...
\end{verbatim}

The flow shell is regenerated with the original decorator structure
but with empty bodies — the populator does not capture statement-level
imperative code, so the generator produces a runnable scaffold and
emits a \texttt{TODO} marker per step:

\begin{verbatim}
class EmailAutoResponderFlow(Flow[EmailAutoResponderFlowState]):
    @start("wait_next_run")
    def fetch_new_emails(self):
        result = EmailFilterCrew().crew().kickoff()
        return result

    @listen(fetch_new_emails)
    def generate_draft_responses(self):
        pass  # TODO: implement step logic
\end{verbatim}

The two ABox-level losses visible in the metrics for this example are
exactly the two losses visible in the source: the YAML quoting style
flips from block scalars (\texttt{>}) to flow scalars, which
canonicalises whitespace inside the literals (the literal-overlap
score for this system is 0.96, not 1.00); and the imperative bodies
of \texttt{fetch\_new\_emails} and \texttt{generate\_draft\_responses}
are not modelled in the A-Box and consequently disappear, which
contributes the bulk of the 0.211-point gap to perfect
\textsc{ast}~F$_1$. Both losses are by design: the ontology represents
the configurable surface of an agentic system, not its runtime
implementation.

\subsection{Cross-framework round-trip}\label{sec:cross-framework}

The cross-framework round-trip extracts a system in source framework
$f$ and re-generates it in a different target framework $t$, before
extracting again. Here the AST diff is dropped (the two sides speak
different languages) and replaced by the
\texttt{mapping\_conformance} metric, which counts canonical
construct-pair occurrences between CrewAI~Flow and LangGraph along
the rule set in
\texttt{evaluation/mappings/crewai\_flow\_to\_langgraph.yml} (graph
class, node-to-decorated, entry point, sequential edge, conditional,
state reducer, kickoff, fan-in/and, fan-in/or). The two cross-fw
directions exercise different parts of the pipeline: CrewAI~$\to$~LangGraph
stresses the LangGraph \emph{generator}; LangGraph~$\to$~CrewAI stresses
the CrewAI generator and the ability of the CrewAI parser to recover
a Flow-shaped semantics from a LangGraph-shaped graph.

\begin{table}[ht]
  \centering
  \small
  \caption{Cross-framework round-trip results.
           CF~$\to$~LG transforms the seven CrewAI examples through
           LangGraph and back; LG~$\to$~CF transforms the six
           LangGraph examples through CrewAI and back.}
  \label{tab:roundtrip-cross}
  \resizebox{\textwidth}{!}{%
  \begin{tabular}{l l c c c c c}
    \toprule
    Direction & Example & Triple F$_1$ & Property F$_1$ & Individual F$_1$ & Fuzzy avg. & Mapping\\
    \midrule
    CF~$\to$~LG & academic-research-flow & 0.131 & 0.833 & 0.688 & 0.873 & 0.775\\
    CF~$\to$~LG & code-review            & 0.148 & 0.902 & 0.796 & 0.829 & 0.687\\
    CF~$\to$~LG & comprehensive          & 0.147 & 0.796 & 0.736 & 0.865 & 0.793\\
    CF~$\to$~LG & content-pipeline       & 0.095 & 0.923 & 0.806 & 0.804 & 0.695\\
    CF~$\to$~LG & email-flow             & 0.100 & 0.900 & 0.806 & 0.821 & 0.750\\
    CF~$\to$~LG & self-eval-loop-flow    & 0.144 & 0.872 & 0.709 & 0.821 & 0.674\\
    CF~$\to$~LG & tech-blog              & 0.303 & 0.899 & 0.746 & 0.883 & 0.167\\
    \addlinespace
    LG~$\to$~CF & ReAct                  & 0.578 & 0.970 & 0.923 & 1.000 & 0.238\\
    LG~$\to$~CF & drafter                & 0.588 & 1.000 & 0.944 & 1.000 & 0.259\\
    LG~$\to$~CF & memoryagent            & 0.586 & 1.000 & 0.944 & 1.000 & 0.583\\
    LG~$\to$~CF & ragagent               & 0.658 & 0.959 & 0.943 & 1.000 & 0.444\\
    LG~$\to$~CF & research-assistant     & 0.653 & 0.946 & 0.943 & 1.000 & 0.522\\
    LG~$\to$~CF & tech-blog              & 0.674 & 1.000 & 0.966 & 1.000 & 0.561\\
    \midrule
    \multicolumn{2}{l}{Mean (CF~$\to$~LG)} & 0.153 & 0.875 & 0.755 & 0.842 & 0.649\\
    \multicolumn{2}{l}{Mean (LG~$\to$~CF)} & 0.623 & 0.979 & 0.944 & 1.000 & 0.435\\
    \bottomrule
  \end{tabular}
  }
\end{table}

The cross-framework numbers are interpreted differently from the
same-framework numbers, because the two A-Boxes $G_1, G_2$ in
Table~\ref{tab:roundtrip-cross} now describe systems implemented in
\emph{different} frameworks. Three findings stand out.

\textbf{Triple $F_1$ collapses but property/individual $F_1$ remain
high.} For CF~$\to$~LG the mean triple $F_1$ falls to 0.153 while
property $F_1$ stays at 0.875 and individual $F_1$ at 0.755. This
spread is the canonical fingerprint of a controlled-translation
round-trip: the entities and the ontology vocabulary survive (a
LangGraph extraction of a generated LangGraph crew still recognises
the same agents, tools and teams), but their concrete triples differ
because each framework's parser emits framework-specific predicates
(e.g., \texttt{taskAssignedToAgent} in CrewAI vs.\
\texttt{nextStep} chains in LangGraph) that the other side does not
produce. The fuzzy-match score (0.842 for CF~$\to$~LG, 1.000 for
LG~$\to$~CF) confirms that the misalignment is largely at the
\emph{predicate} level, not at the entity level: every named entity
in $G_1$ has a high-similarity counterpart in $G_2$.

\textbf{The two directions are asymmetric.} LG~$\to$~CF preserves
roughly four times more triples than CF~$\to$~LG (0.623 vs.\ 0.153).
The asymmetry reflects an information-theoretic gap between the
frameworks: CrewAI Flows encode imperative event-driven control with
\texttt{@start}/\texttt{@listen}/\texttt{@router} decorators that
serialise into rich \texttt{Orchestration} subgraphs in the A-Box;
when these are mapped onto a LangGraph \texttt{StateGraph} the
imperative semantics is flattened into edge declarations and
re-extracted under a smaller predicate set, so $G_2$ has fewer
triples than $G_1$. The reverse direction is more lossless because
LangGraph's vocabulary is a near-subset of the constructs that the
CrewAI parser can reconstruct.

\textbf{Mapping conformance highlights remaining gaps in the
generators.} The canonical-rule conformance averages 0.649 for
CF~$\to$~LG and 0.435 for LG~$\to$~CF, with substantial per-example
variance. Inspection of the per-rule scores shows that the
\texttt{graph\_class}, \texttt{kickoff} and \texttt{entry\_point}
rules are matched almost everywhere, while \texttt{state\_reducer}
(\texttt{Annotated[\,T,\,reducer]} $\leftrightarrow$ Pydantic
\texttt{BaseModel} default) and the conditional/fan-in rules drive
the remaining loss: the LangGraph generator currently emits a
\texttt{TypedDict} state schema without reducer annotations, and the
CrewAI generator does not yet round-trip
\texttt{add\_conditional\_edges} into a \texttt{@router}-decorated
method. These two gaps are the most actionable items surfaced by the
benchmark and are listed as future work in Chapter~9.

% =====================================================================
\section{Gold-Anchored Comparison with LLM Baseline}\label{sec:llm-baseline}
% =====================================================================

To situate the deterministic AST-based pipeline against the LLM-driven
approach taken by AgentO~\cite{EKEK26}, the harness includes an
\texttt{extract-llm} command that reuses the AgentO extraction prompt
verbatim, instantiated with the \textsc{AgentOscin} schema in place of
AgentO's (prompt template: \texttt{oscin/prompts/llm\_extraction.md}).
Since \textsc{AgentOscin} is a conservative extension of AgentO
(Chapter~5), the prompt is directly transferable: the LLM is given the
extended schema and the raw source files in a single call and asked to
emit populated Turtle. The comparison therefore evaluates the AgentO
extraction \emph{method} on the \textsc{AgentOscin} ontology, against
the same gold A-Boxes and the same metrics used for the deterministic
pipeline.

\paragraph{Gold reference.} For six systems chosen to span all three
frameworks and the architectural patterns covered by the corpus —
\texttt{crewai/email-flow} (event-driven flow with sub-crew),
\texttt{crewai/academic-research-flow} (hierarchical crew with
human-in-the-loop and conditional review),
\texttt{langgraph/ReAct} (manual \texttt{StateGraph} reasoning loop),
\texttt{langgraph/research-assistant} (router-based retrieval graph),
\texttt{autogen/code-review} and \texttt{autogen/content-pipeline}
(round-robin group chats) — we hand-authored a gold A-Box from the
source files following the annotation guideline in
\texttt{docs/annotation\_guideline.md}. The six gold graphs together
contain 505 ABox triples over 121 individuals using 41 distinct
ontology properties.

\paragraph{Setup.} We ran the LLM extractor with provider
\texttt{openai}, model \texttt{gpt-4o}, \texttt{max\_tokens}=16\,384
and the default temperature, once per system. One of the six outputs
needed minor TTL repair (a missing end-of-statement \texttt{.} after a
literal); LLM outputs that mix schema and individuals into the
ontology namespace are re-based into a separate instance namespace
before scoring (the deterministic AST pipeline emits the two
namespaces correctly without post-processing). Calls cost approximately
USD~0.05 per extraction at the prompt sizes in this corpus.

\paragraph{Metrics.} Individual and property $F_1$ are computed as in
the round-trip evaluation. We report two triple-level scores. The
\emph{exact} triple $F_1$ matches $(s,p,o)$ tuples after URI
local-name normalisation, and is sensitive to the subject naming
scheme. The \emph{aligned} triple $F_1$ first computes a greedy
bipartite alignment of candidate individuals to gold individuals,
restricted to pairs that share at least one \texttt{rdf:type} and
scored by Jaccard overlap of (predicate,literal) and
(predicate,object-class) features, then rewrites the candidate's
individual URIs through the alignment before recomputing the triple
set. This metric isolates content fidelity from the choice of local
names. Literal overlap reports recall of the gold's literal values
(prompt strings, descriptions, identifiers) regardless of which
subject they attach to.

\begin{table}[ht]
  \centering
  \footnotesize
  \caption{Gold-anchored extraction metrics on the 6-system subset.
    Each system contributes three rows: the gold reference (sizes
    only), the deterministic AST extractor and the LLM extractor.
    \emph{Indiv.}, \emph{ABox} and \emph{Prop.} are intrinsic A-Box
    sizes (named individuals, ABox triples, distinct ontology
    properties). $F_1$ scores are computed against the gold reference
    of the same row group; \emph{aligned} triple $F_1$ uses the
    class-aware bipartite alignment described in the text;
    \emph{lit.\ ovr.} is the gold-literal recall.}
  \label{tab:gold-vs-extractors}
  \resizebox{\textwidth}{!}{%
  \begin{tabular}{l l l r r r c c c c c}
    \toprule
    Framework & Example & Source
              & Indiv. & ABox & Prop.
              & Ind.\ $F_1$ & Prop.\ $F_1$ & Triple $F_1$
              & Aligned $F_1$ & Lit.\ ovr.\\
    \midrule
    CrewAI    & email-flow              & gold          & 35 & 148 & 34 & --    & --    & --    & --    & --    \\
              &                         & deterministic & 39 & 187 & 44 & 0.946 & 0.872 & 0.322 & 0.812 & 0.933 \\
              &                         & LLM (gpt-4o)  & 20 & \phantom{0}75 & 19 & 0.727 & 0.566 & 0.000 & 0.457 & 0.311 \\
    \addlinespace
    CrewAI    & academic-research-flow  & gold          & 33 & 143 & 41 & --    & --    & --    & --    & --    \\
              &                         & deterministic & 45 & 188 & 54 & 0.821 & 0.821 & 0.308 & 0.725 & 0.844 \\
              &                         & LLM (gpt-4o)  & 23 & \phantom{0}79 & 23 & 0.571 & 0.469 & 0.000 & 0.315 & 0.289 \\
    \addlinespace
    LangGraph & ReAct                   & gold          & 14 & \phantom{0}56 & 23 & --    & --    & --    & --    & --    \\
              &                         & deterministic & 20 & \phantom{0}88 & 34 & 0.765 & 0.702 & 0.236 & 0.500 & 0.650 \\
              &                         & LLM (gpt-4o)  & 17 & \phantom{0}51 & 21 & 0.581 & 0.318 & 0.000 & 0.280 & 0.400 \\
    \addlinespace
    LangGraph & research-assistant      & gold          & 16 & \phantom{0}63 & 26 & --    & --    & --    & --    & --    \\
              &                         & deterministic & 27 & 127 & 39 & 0.698 & 0.677 & 0.168 & 0.463 & 0.636 \\
              &                         & LLM (gpt-4o)  & 20 & \phantom{0}76 & 20 & 0.667 & 0.391 & 0.000 & 0.273 & 0.182 \\
    \addlinespace
    AutoGen   & code-review             & gold          & 11 & \phantom{0}45 & 16 & --    & --    & --    & --    & --    \\
              &                         & deterministic & 26 & 110 & 29 & 0.595 & 0.711 & 0.219 & 0.503 & 0.733 \\
              &                         & LLM (gpt-4o)  & 23 & \phantom{0}83 & 22 & 0.529 & 0.421 & 0.000 & 0.344 & 0.533 \\
    \addlinespace
    AutoGen   & content-pipeline        & gold          & 12 & \phantom{0}50 & 16 & --    & --    & --    & --    & --    \\
              &                         & deterministic & 27 & 116 & 29 & 0.615 & 0.711 & 0.265 & 0.518 & 0.778 \\
              &                         & LLM (gpt-4o)  & 14 & \phantom{0}54 & 17 & 0.846 & 0.545 & 0.000 & 0.500 & 0.389 \\
    \bottomrule
  \end{tabular}
  }
\end{table}

\begin{table}[ht]
  \centering
  \small
  \caption{Aggregate gold-anchored $F_1$ (mean across the six
    systems in Table~\ref{tab:gold-vs-extractors}).}
  \label{tab:gold-aggregate}
  \begin{tabular}{l c c c c c}
    \toprule
    Source & Ind.\ $F_1$ & Prop.\ $F_1$ & Triple $F_1$ (exact)
           & Triple $F_1$ (aligned) & Lit.\ overlap\\
    \midrule
    Deterministic AST     & 0.740 & 0.749 & 0.253 & 0.587 & 0.762\\
    LLM (gpt-4o)          & 0.653 & 0.452 & 0.000 & 0.361 & 0.351\\
    \bottomrule
  \end{tabular}
\end{table}

The numbers in
Tables~\ref{tab:gold-vs-extractors}--\ref{tab:gold-aggregate} support
four conclusions.

\textbf{Coverage gap.} The LLM produces between one third (CrewAI
\texttt{email-flow}: $75/148$ relative to gold) and slightly above
gold size (AutoGen \texttt{content-pipeline}: $54/50$) of the gold
triple count, while the deterministic pipeline systematically
\emph{over-produces} (1.27$\times$--2.44$\times$ gold size) because
its parser materialises orchestration nodes (every \texttt{@start}/
\texttt{@listen} edge, every \texttt{StateGraph} edge, every
group-chat seat) that the gold leaves implicit. The literal-overlap
column makes the coverage difference concrete: 0.762 of the gold's
identifying strings (prompt instructions, tool names, descriptions)
are present somewhere in the deterministic A-Box, against 0.351 for
the LLM.

\textbf{Property and individual recall.} Against gold, the
deterministic pipeline averages 0.749 property $F_1$ and 0.740
individual $F_1$; the LLM averages 0.452 and 0.653. Individual $F_1$
is comparatively close because both extractors recover most agent and
tool entities; the deterministic pipeline's larger lead on property
$F_1$ comes from the orchestration vocabulary
(\texttt{hasStartStep}, \texttt{listensToEvent},
\texttt{hasNextStep}, \texttt{hasTerminationCondition}) that the LLM
prompt does not reliably surface.

\textbf{Exact triple $F_1$ collapses to zero for the LLM.} Because
triples are scored as $(\text{subject},\text{predicate},\text{object})$
tuples and the LLM chooses different individual names from the gold,
\emph{no} normalised triple in any LLM output matches gold on any of
the six systems. The deterministic pipeline retains a non-trivial
exact triple $F_1$ of 0.253 because it inherits a deterministic naming
convention (\texttt{:Agent\_<source\_id>}, \texttt{:Tool\_<name>}) that
the gold annotator was instructed to follow. The aligned metric
removes this naming-scheme effect: under content alignment the
deterministic pipeline scores 0.587 and the LLM 0.361, a gap of 0.23
absolute that reflects \emph{actual} content disagreement rather than
URI choice.

\textbf{The LLM is non-stationary even on a single run.} Even with the
same model, prompt and parameters the LLM omits different ontology
slots on different systems (property $F_1$ ranges 0.32--0.57 across
the six gold-anchored runs). The deterministic pipeline's main
practical advantage over the LLM is therefore not only that it
recovers more content per run, but that the triples it produces are
\emph{stable}: every re-run yields a byte-identical A-Box, and the
gold-anchored scores in Table~\ref{tab:gold-vs-extractors} are
deterministic functions of the source files. For a pipeline whose
downstream artefacts (generated source, mapping conformance,
CQ-survivability checks) depend on the A-Box being reproducibly
named, this matters more than the absolute $F_1$ numbers.

The point of this comparison is therefore not that one approach
dominates the other, but to characterise the trade-off precisely:
the deterministic pipeline is reproducible and its triples are stable
across runs (every triple is produced by a documented mapping rule
applicable to every input), but it is bounded by the patterns its
parser knows; the LLM baseline is unbounded in pattern coverage — it
can in principle map a free-form Python module to ontology
individuals — but its triple-level output is non-stationary and would
need a calibration layer (e.g., deterministic naming, schema-driven
validation) before it could participate in the round-trip.

% =====================================================================
\section{Discussion of Results}\label{sec:transformation-discussion}
% =====================================================================

We organise the discussion around the three transformation-related
research questions stated in Chapter~3, since each table above
addresses a different one.

\paragraph{RQ-2 (extraction).}
The intrinsic statistics in Table~\ref{tab:extraction-intrinsic}
together with the gold-anchored scores in
Table~\ref{tab:gold-vs-extractors} answer the operational form of
\textit{RQ-2}. On the full 21-system corpus the deterministic
parser--populator chain produces non-trivial, ontology-broad A-Boxes
for every system (553--678 triples, 18--45 individuals, 26--56
properties used). On the 6-system subset for which a hand-authored
gold A-Box exists, the same pipeline attains 0.749 property $F_1$,
0.740 individual $F_1$ and 0.587 aligned triple $F_1$ against gold
(Table~\ref{tab:gold-aggregate}). The LLM baseline on the same gold
attains 0.452 property $F_1$, 0.653 individual $F_1$ and 0.361
aligned triple $F_1$, and collapses to 0.000 on the exact triple
metric because it does not commit to a deterministic naming scheme.
The deterministic pipeline therefore leads by 0.23--0.30 absolute on
every level scored against gold, and is additionally byte-stable
across re-runs. We answer RQ-2 in the affirmative for the
deterministic pipeline, and characterise the LLM baseline as a
complement (coverage extension) rather than a substitute
(Section~\ref{sec:llm-baseline}).

\paragraph{RQ-3 (generation and round-trip).}
Tables~\ref{tab:generation-validity}, \ref{tab:roundtrip-same}
and~\ref{tab:roundtrip-cross} jointly answer \textit{RQ-3}: an
\textsc{AgentOscin} A-Box contains enough information to regenerate
a runnable target-framework artefact (100\,\% syntax rate, 1.000
import rate for two of three targets) and to round-trip the system
within the same framework with high semantic fidelity (mean triple
$F_1$ 0.847, property/individual $F_1\geq 0.97$). Cross-framework
the same metric set splits into two regimes: the entity inventory
and the property vocabulary survive almost unchanged (property
$F_1\geq 0.88$, fuzzy-alignment $\geq 0.84$), but the triple-level
overlap drops to 0.15--0.62 because the predicates the two parsers
emit are framework-specific. The walked-through example
(Section~\ref{sec:walkthrough}) shows where the residual gap comes
from concretely: literal-value paraphrasing in the YAML emitter and
the deliberate exclusion of statement-level imperative code from the
ontology surface. The CQ-survivability check
(Table~\ref{tab:cq-survivability}) strengthens the answer: the
questions the ontology is designed to answer continue to return the
same answer multisets on the round-tripped graph for $\geq 0.78$ of
systems on every CQ, and never drop below a Jaccard of 0.83 even
where the multisets differ.

\paragraph{RQ-4 (evaluation methods).}
The seven-metric harness in Table~\ref{tab:metric-catalogue}, applied
across three pipelines, the LLM baseline and the CQ-survivability
check, is itself the answer to \textit{RQ-4}: each metric isolates a
different failure mode (intrinsic — empty extraction; pairwise — A-Box
drift; fuzzy — entity rename; AST — syntactic regression; syntax/import
— validity floor; mapping — semantic loss across frameworks;
SPARQL — query-level survival), and together they let us attribute
every drop in any aggregate score in this chapter to a localised cause
in either the populator or one of the generators. The CrewAI 0.721
import-rate, the LangGraph 0.659 \textsc{ast}~F$_1$ and the LG~$\to$~CF
0.435 mapping conformance are three examples where the harness
produces an actionable diagnosis (Chapter~9). The methodology is
therefore \emph{semantically} round-trip-stable, while
\emph{syntactic} round-trip stability remains framework-bounded.

The benchmark also surfaces three concrete limitations that motivate
future work. First, the CrewAI generator's reliance on external tool
references — visible as the 0.721 mean import rate in
Table~\ref{tab:generation-validity} — could be replaced by an
optional stub-tool emitter that produces a runnable
\texttt{BaseTool} subclass for every external reference; this would
trade an explicit modelling decision (preserve foreign tool
identity) for runnability. Second, the LangGraph generator does not
yet emit \texttt{Annotated[\,T,\,reducer]} state-field
annotations, which depresses both the \textsc{ast}~F$_1$ in the
same-framework round-trip (LangGraph mean 0.659) and the
\texttt{state\_reducer} mapping rule in the cross-framework
round-trip; adding reducer support is a localised change in
\texttt{oscin/generators/langgraph\_generator.py}. Third, neither
generator currently round-trips \texttt{ConditionalStep} into a
fully-formed routing function, which is the dominant contributor to
the 0.435 mapping conformance mean for LG~$\to$~CF; lifting this would
push that direction towards parity with CF~$\to$~LG.

Finally, the corpus is a convenience sample of public examples and is
small enough that per-example variance dominates per-framework
variance for some metrics (notably the 0.167 mapping conformance for
CrewAI \texttt{tech-blog}~$\to$~LangGraph, which is an outlier driven
by a flat one-step crew that has no LangGraph-side construct to
match). Threats to the validity of the numbers, including this
sample-size concern, are revisited in Chapter~8.
